\section{Fixed-parameter forest selection beyond two-row columns}
\label{sec:exceptional-forest}
\label{sec:exceptional}

The two-row theorem of \cite{IncomingDualNormalization} reduces a restricted
dual-forest choice to directed reachability and independent representatives
in a graphic matroid.  The remaining original native column shape is
$(1,1,1)$: deleting two of its rows allows the third to peel.  Such a rule
has two premises and cannot be represented by ordinary directed
reachability alone.  This section gives an exact reduction when the rows
participating in these exceptional rules form a small interface.  The
parameter is independent of the number of interior dual vertices.

\subsection{The finite problem and its parameter}
\label{subsec:exceptional-problem}

Let $D$ be a connected finite undirected multigraph, with a distinguished
root and vertex set $V$.  Let $R$ be a finite row set.  Each row $r\in R$
labels either an edge of $D$ or an ineligible record that cannot be selected
in a forest.  Eligible row records are in bijection with the edges of
$D$, so distinct parallel edges still have distinct row labels.
Loops and parallel edges are allowed.  Write
\[
                    A=(a_{rc})\in\mathbb{N}^{R\times C}
\]
for a nonnegative integral matrix, supplied by sparse occurrence lists.
In the geometric application, a row is an attaching two-handle and a
column is a retained one-handle cocore.  That interpretation requires the
separate source-authentication and geometric cancellation theorems; the
finite operator below does not establish it from the matrix alone.

A spanning tree $T\subseteq R$ deletes its rows.  A \emph{unit-column peel}
on a current live row set $L$ removes a row $r\in L$ when some column $c$
satisfies
\[
                      \sum_{s\in L}a_{sc}=a_{rc}=1.
\]
Its column may also be deleted, since it has no other live occurrence.
The decision problem is whether some spanning tree leaves a row restriction
that can be completely peeled.  By the raw-cancellation theorem of
\cite{IncomingDualNormalization}, this is also the exact question for
complete positive raw cancellation in the specified occurrence model.
A negative answer concerns that model and its fixed geometric charts.

Call a column \emph{active} if it has an entry equal to one.  An inactive
column cannot supply a peel after any row deletion, so it can be ignored
for this decision.  For a column $c$, let
\[
         W_c=\{r\in R:a_{rc}>0\},\qquad
         E=\{c\in C:c\text{ is active and }|W_c|>2\}.
\]
The \emph{exceptional interface} and its size are
\begin{equation}
          U=\bigcup_{c\in E}W_c,\qquad u=|U|.
          \label{eq:exceptional-interface}
\end{equation}
Inert columns with large supports do not contribute to $U$.  Conversely,
the support of an exceptional active column contributes \emph{all} its
rows, including rows whose own entry in that column exceeds one.

Throughout this section, $N$ denotes the encoded input size, including
graph records, sparse incidences, and the binary encodings of counts.
Counts can be capped at two for the zero/one decisions used here.  Such
capping must not be used to justify a separately asserted parity property
of the original geometric chain complex.

\subsection{Baseline reachability and mandatory source components}
\label{subsec:exceptional-baseline}

Temporarily omit the columns of $E$.  Build a directed graph $P$ on $R$
from the remaining active columns.  A singleton column with entry one at
$r$ makes $r$ an automatic seed; let $Z$ be the set of these seeds.  A
column supported on $\{r,s\}$ gives the arc $s\to r$ if $a_{rc}=1$, and
the arc $r\to s$ if $a_{sc}=1$.  No restriction is imposed on the other
positive entry.  Reachability includes paths of length zero.

\begin{lemma}[Exact baseline activation]
\label{lem:exceptional-baseline}
If a row set $Q$ is deleted initially, the rows deleted or peeled using
only the baseline columns are exactly
\[
                         \operatorname{Reach}_P(Z\cup Q).
\]
For an exceptional column $c$, its additional exact activation rules are
\begin{equation}
             W_c\setminus\{r\}\ \Longrightarrow\ r
             \qquad(r\in W_c,\ a_{rc}=1).
             \label{eq:exceptional-horn-rules}
\end{equation}
Consequently closure under baseline reachability and these rules is
exactly the set of initially deleted or subsequently peelable rows.
\end{lemma}

\begin{proof}
An automatic seed can peel without another row deletion.  For a two-row
column, deleting the other supported row makes its unit entry a legal
pivot.  Thus every directed path gives a valid sequence of possible
activations.  Conversely, each baseline peel is either an automatic seed
or follows an arc whose tail is already activated.  Induction on a peel
trace proves the first assertion, recovering the reachability model of
\cite{IncomingDualNormalization}.

For an exceptional column, a live row $r$ can peel in that column exactly
when its entry is one and every other supported row has been deleted.
These are precisely the premises of \eqref{eq:exceptional-horn-rules}.
Interleaving these rules with baseline activations proves exactness for
the full matrix.  Continuing to retain a column in this logical
description after its pivot is harmless: at that pivot it had no other
live owner, and therefore cannot later remove a different live row.
\end{proof}

Let $S_1,\ldots,S_q$ be the source strongly connected components of $P$
that contain no automatic seed and are disjoint from $U$.  They are
pairwise disjoint, and their number need not be small.  Put
\begin{equation}
       M=\bigcup_{i=1}^{q}S_i,\qquad
       B=\operatorname{Reach}_P(Z\cup M),\qquad B_U=B\cap U.
       \label{eq:exceptional-base}
\end{equation}
The definition of $B$ treats all rows of $M$ as activated for convenience;
the algorithm will actually require only one selected forest row from
each $S_i$.

\begin{lemma}[Mandatory groups and an interface completion criterion]
\label{lem:exceptional-obligations}
Every initially deleted row set that permits full peeling must meet each
$S_i$.  A single initially deleted row in $S_i$ activates the entire
component by baseline reachability.  If $Q$ meets all these groups, then
full activation of $R$ is equivalent to activation of every row in $U$.
\end{lemma}

\begin{proof}
The component $S_i$ has no incoming arc in the condensation graph, no
automatic seed, and no row in the support of an exceptional rule.  Hence
neither baseline propagation from outside nor an exceptional activation
can create its first activated row.  It must contain an initially deleted
row.  Strong connectivity then activates all its rows and every baseline
descendant.

For the last assertion, activation of all of $R$ plainly activates $U$.
Conversely, every source strongly connected component either contains an
automatic seed, is one of the groups $S_i$, or intersects $U$.  In the
third case activation of all of $U$ supplies a row of that component;
in the first two cases a row is already supplied by $Z$ or $Q$.
Every source component is therefore activated in full.  Each component
of a finite directed acyclic condensation graph is reachable from a
source, so all of $R$ becomes activated.
\end{proof}

\subsection{Compressing all nonlinear interaction to the interface}
\label{subsec:exceptional-compression}

For any row $r\in R$, define its reachability signature
\begin{equation}
                   \sigma(r)=\operatorname{Reach}_P(\{r\})\cap U.
                   \label{eq:exceptional-signature}
\end{equation}
On subsets of $U$, define $\operatorname{Hcl}$ to be the least closure
under the following finite rules:
\begin{enumerate}
\item $v\Longrightarrow w$ whenever $v,w\in U$ and $w$ is reachable
      from $v$ in $P$;
\item every implication in \eqref{eq:exceptional-horn-rules}.
\end{enumerate}
All premises and conclusions in the second family lie in $U$ by its
definition.  Thus this closure can be computed without retaining the
non-interface vertices as dynamic states.

\begin{lemma}[Exact compressed Horn closure]
\label{lem:exceptional-compressed}
Let $Q\subseteq R$ meet every mandatory group $S_i$.  The interface rows
eventually activated from the initial deletion of $Q$ are exactly
\begin{equation}
  \operatorname{Hcl}\left(
        B_U\ \cup\!\bigcup_{r\in Q\setminus M}\sigma(r)
                      \right).
  \label{eq:exceptional-compressed-closure}
\end{equation}
In particular, $A[R\setminus Q,C]$ is completely peelable if and only if
the set in \eqref{eq:exceptional-compressed-closure} is $U$.
\end{lemma}

\begin{proof}
Each mandatory group has a selected row, so all of $B$ can be activated
by the baseline rules.  Every $\sigma(r)$ in the displayed expression
is reachable from an actual initially deleted row.  Every interface
reachability rule abbreviates a valid baseline path, and every exceptional
rule is exact by Lemma~\ref{lem:exceptional-baseline}.  The displayed
closure is therefore contained in the actual activated interface.

For the reverse inclusion, consider an activation trace and separate its
baseline propagation from its exceptional rule firings.  Before any
exceptional firing, each activated interface row is reached from $Z$ or
$Q$.  Paths from $Z$ or $Q\cap M$ are covered by $B_U$; paths from the
other rows of $Q$ are covered by their signatures.  If an exceptional rule
fires, its head belongs to $U$, and all subsequent baseline propagation
from this head to interface rows is represented by interface reachability
rules.  The exceptional implication itself is retained without any
change of premises.  Induction on the exceptional firings proves that
every actual interface activation occurs in
\eqref{eq:exceptional-compressed-closure}.  The final equivalence follows
from Lemma~\ref{lem:exceptional-obligations}.
\end{proof}

The compression is algorithmic as well as logical.  Reverse reachability
from each $v\in U$ computes all signatures: a row $r$ receives the target
$v$ exactly when $r$ reaches $v$ in $P$.  Forward reachability from the
interface rows computes the interface-to-interface rules.  These
computations use a polynomial number of graph operations.  Repeatedly
scanning the finite rule list until no target is added computes
$\operatorname{Hcl}$; there are at most $u$ strict additions.

\subsection{Partial partitions and independent representatives}
\label{subsec:exceptional-partitions}

A \emph{partial partition} of $U$ consists of a subset $U_0\subseteq U$
and a partition
\[
                   \Pi=\{B_1,\ldots,B_\ell\}
\]
of $U_0$ into nonempty blocks.  The empty subset has the empty partition.
The interpretation is that $U_0$ is a guessed collection of interface
targets activated directly by selected tree rows through baseline paths;
one additional tree row is to supply each block.  This is a witness
assignment, not an assertion that all possible tree seeds have already
been identified.

For each partial partition, first test
\begin{equation}
                    \operatorname{Hcl}(B_U\cup U_0)=U.
                    \label{eq:exceptional-partition-test}
\end{equation}
If this fails, discard the partition.  Let $R_{\mathrm{el}}$ consist of
the rows labelling eligible nonloops of $D$.  Otherwise define the
candidate group for its $j$th block by
\begin{equation}
       G_j=\{r\in R_{\mathrm{el}}\setminus M:
                                      B_j\subseteq\sigma(r)\}.
       \label{eq:exceptional-block-group}
\end{equation}
Seek distinct representatives, one from each member of the group list
\begin{equation}
                    S_1,\ldots,S_q,G_1,\ldots,G_\ell,
                    \label{eq:exceptional-group-list}
\end{equation}
whose graph edges form a forest.  Ineligible rows and loops can also be
removed from the mandatory candidate groups when performing this test.
An empty group makes that particular test fail.

\begin{lemma}[Overlapping groups reduce to graphic/partition intersection]
\label{lem:exceptional-clones}
The independent-representatives problem for
\eqref{eq:exceptional-group-list} is decidable in polynomial time, even
though the block groups $G_j$ may overlap.  It requires at most
$(\ell+1)|R|$ graph-edge copies.
\end{lemma}

\begin{proof}
Make one edge copy $(r,H)$ for each permitted membership of row $r$ in a
group $H$ from \eqref{eq:exceptional-group-list}.  Give this copy the same
endpoints as $r$ and color it by $H$.  Copies of a single original edge
are parallel in the resulting multigraph.  Two such copies cannot both
belong to a forest, since they form a two-edge cycle.  Graphic
independence therefore enforces both distinct original representatives
and independence of the original edge set.  The partition matroid that
allows at most one copy of each color enforces the group constraints.
A common independent set of size $q+\ell$ exists exactly when the desired
representatives exist.

This is the standard graphic/partition matroid-intersection problem,
with polynomial independence tests and a polynomial augmenting-path
algorithm~\cite{Edmonds1970,Rado1942}.  It is also Rado's independent-transversal formulation,
already used in \cite{IncomingDualNormalization}; no new
matroid-intersection theorem is needed.  The mandatory groups are
disjoint and contribute at most $|R|$ memberships in total.  Every other
row contributes at most one copy per block group, hence at most
$\ell|R|$ additional memberships.  This proves the stated size bound.
\end{proof}

The exclusion of $M$ from the block candidates loses no useful initial
activation.  Once the mandatory groups have their representatives, all
signatures of rows in $M$ are already included in the baseline contribution
$B_U$.  The overlapping candidate groups outside $M$ must nevertheless be
handled explicitly; treating them as disjoint colors on the original
edge set would be incorrect.

\begin{theorem}[Exact bounded-interface forest selection]
\label{thm:exceptional-selector}
The following finite procedure decides whether some spanning tree of $D$
leaves $A$ completely peelable:
\begin{enumerate}
\item Construct the baseline graph, automatic seeds, mandatory groups,
      signatures, and interface closure.
\item Enumerate all partial partitions of $U$.  Retain only those
      satisfying \eqref{eq:exceptional-partition-test}.
\item For each retained partition, solve the independent-representatives
      problem \eqref{eq:exceptional-group-list}.
\item If representatives form a forest $Q$, extend $Q$ to a spanning
      tree $T$ of $D$, recompute a complete peel order on
      $A[R\setminus T,C]$, and return $T$ and that order.
\item If all partial partitions have failed, report that no spanning
      tree in the specified finite class permits full peeling.
\end{enumerate}
The algorithm makes at most $\operatorname{Bell}(u+1)$ polynomial
independent-representatives calls.  Its bit complexity is
\begin{equation}
       2^{O(u\log(u+2))}\operatorname{poly}(N).
       \label{eq:exceptional-fpt}
\end{equation}
\end{theorem}

\begin{proof}
Suppose the algorithm finds representatives forming $Q$.  Each mandatory
group is met.  Each block $B_j$ is contained in the actual signature of
its representative, so the actual initial interface activation contains
$B_U\cup U_0$.  Equation~\eqref{eq:exceptional-partition-test} and
Lemma~\ref{lem:exceptional-compressed} show that deleting $Q$ permits
activation of every row.  Thus $A[R\setminus Q,C]$ is completely peelable.
Every forest in a connected graph extends to a spanning tree.  Deleting
additional rows preserves peelability: restrict an existing peel order
to the rows still present; at each remaining pivot its own occurrence is
one and the other live occurrences remain absent.  Hence the extension
$T$ also succeeds.  Recomputing the final trace provides an ordinary
tree-and-peeling certificate on the actual chosen complement.  This
proves soundness.

For completeness, suppose a successful spanning tree $T$ exists.
Lemma~\ref{lem:exceptional-obligations} shows that it meets each $S_i$;
choose one of its rows from each group.  Define
\[
                   U_0=\bigcup_{r\in T\setminus M}\sigma(r).
\]
Since $T$ succeeds, Lemma~\ref{lem:exceptional-compressed} implies
\eqref{eq:exceptional-partition-test}.  For each $v\in U_0$, choose an
actual witnessing row $r(v)\in T\setminus M$ with
$v\in\sigma(r(v))$.  Partition $U_0$ by equality of these witnesses.
The resulting nonempty blocks have distinct representatives in
$T\setminus M$, and each block is contained in the signature of its
representative.  Together with the chosen mandatory representatives,
they form a subset of $T$ and are therefore independent.  The algorithm
enumerates this partial partition, and its representatives test succeeds.
This argument proves why at most $u$ extra representatives suffice;
it does not assume that arbitrary tree edges can be replaced by source
rows while preserving independence.

For the count, adjoin one distinguished element $\ast$ to $U$.  A
partition of $U\cup\{\ast\}$ corresponds bijectively to a partial
partition of $U$: the block containing $\ast$ records the omitted
elements, and its other blocks are the blocks of the partial partition.
There are exactly $\operatorname{Bell}(u+1)$ objects.  Label the elements
of $U$ by $1,\ldots,u$ and give $\ast$ the label zero.  A partition is
determined by assigning each ordinary element the least label in its
block.  Thus, with the empty-interface case interpreted separately,
\[
                 \operatorname{Bell}(u+1)\leq (u+1)^u.
\]
Every block list has $\ell\leq u$, so
Lemma~\ref{lem:exceptional-clones} gives a graph of polynomial size for
each call.  Candidate membership uses ordinary subset tests on
$u$-element signatures.  Baseline reachability, interface closure,
spanning-tree extension, and the final peel require polynomial time in
the encoded input.  Integer input reading and capping are polynomial in
the binary bit lengths.  Combining these observations yields
\eqref{eq:exceptional-fpt}.
\end{proof}

\begin{remark}[The empty interface and a harmless implementation refinement]
\label{rem:exceptional-zero}
If $u=0$, there is one partial partition, its block list is empty, and
the procedure reduces to the prior polynomial two-row theorem.  More
generally, one may first compute
$U_{\mathrm{base}}=\operatorname{Hcl}(B_U)$ and enumerate partial
partitions only of $U\setminus U_{\mathrm{base}}$.  Targets already
activated by the mandatory groups and automatic seeds need no additional
witness.  Any successful-tree witness assignment can simply omit those
targets, so completeness is preserved.  This can decrease the actual
enumeration substantially; the bound in
Theorem~\ref{thm:exceptional-selector} remains valid without relying on it.
\end{remark}

\subsection{Native specialization and certificate semantics}
\label{subsec:exceptional-native}

\begin{proposition}[Native triple-column specialization]
\label{prop:exceptional-native}
Suppose the matrix is the original source-authenticated native core
occurrence matrix for which each unsaturated column total is at most
three.  Let $k$ be the number of its active columns supported on more
than two rows.  Then each such column has shape $(1,1,1)$ and
\[
                             u\leq 3k.
\]
The finite forest selector therefore has complexity
\begin{equation}
                   2^{O(k\log(k+2))}\operatorname{poly}(N).
                   \label{eq:exceptional-native-fpt}
\end{equation}
\end{proposition}

\begin{proof}
A column supported on at least three rows has at least three positive
occurrences.  The total bound of three forces exactly three supported
rows and an entry of one at each.  Its support contributes at most three
rows to the union defining $U$.  Taking the union over the $k$ exceptional
columns proves $u\leq3k$, and substituting into
\eqref{eq:exceptional-fpt} proves the complexity statement.  The geometric
three-occurrence premise is the prior native cocore theorem of
\cite{IncomingDualNormalization}, not a consequence of the finite matrix
algorithm.
\end{proof}

This specialization applies to the original native source-core matrix and
to row or column deletions that preserve its occurrence interpretation.
Adding annular gluing relations or arbitrary boundary rows can increase
a column's support beyond three.  Such enlarged inputs remain covered by
Theorem~\ref{thm:exceptional-selector} with their actual parameter $u$;
they are not automatically covered by the estimate $u\leq3k$.  Similarly,
transporting additional markings through a chosen dual tree requires the
applicable geometric port-preservation theorem.

The parameterized result places no restriction on $|V|-1$, the number of
tree edges to be chosen.  It gives a local quasi-polynomial regime whenever
the maintained encoded input is polynomial in a diagram parameter $n$
and $u\log(u+2)$ is bounded by a fixed power of $\log n$.  In particular,
$u\log(u+2)=O((\log n)^2)$ would give a local
$n^{O(\log n)}$ bound.  No such interface bound for every state of a
complete unknot-recognition hierarchy is proved here, and the cost of
constructing or branching over those states remains a separate problem.

\begin{remark}[Positive replay, exhaustive failure, and budgets]
\label{rem:exceptional-status}
A positive output contains the final spanning tree and a complete
unit-column peel trace.  Their finite correctness can be checked directly
without replaying the partial-partition search or trusting a
matroid-intersection optimality claim.  The enclosing geometric verifier
must still authenticate the rows, columns, dual endpoints, and topological
meaning of the cancellations.

An exhaustive negative result means that no tree works in this fixed
finite class.  The current implementation does not package the entire
exhaustion as a succinct independent negative certificate.  In particular,
a rank defect for one partial partition is only a rejection of that
partition; every possibility must be excluded before the complete
finite-class negative is justified.  Interrupting the enumeration because
of a work allowance, partition allowance, or implementation
limit must instead return an inconclusive status.  None of these negative
or interrupted outcomes is a certificate that the underlying knot is
nontrivial or that an authenticated core is not a handlebody.
\end{remark}

\subsection{Why triangle-cycle parity does not remove all tree choice}
\label{subsec:exceptional-triangle-witness}

The native parity condition constrains the odd support of each source
column to be an even-degree subgraph of the dual graph.  For a unit
triple, one allowed form is a triangle.  It is tempting to conclude that
such short cycle relations make all spanning trees equivalent for raw
peeling.  The following finite witness rules out that inference from
sparsity and parity alone.

\begin{proposition}[An abstract triangle-column sensitivity witness]
\label{prop:exceptional-triangle-witness}
There is a connected simple graph $D$ on nine vertices with 25 edges and
a Boolean matrix with those edges as rows and 17 columns such that:
\begin{enumerate}
\item each column is supported on exactly the three edges of a triangle
      of $D$;
\item one spanning tree leaves a completely peelable row restriction;
\item another spanning tree leaves a row restriction admitting no
      complete positive raw cancellation.
\end{enumerate}
Thus the three-occurrence bound and dual-cycle parity condition together
do not imply dual-tree invariance for the abstract finite problem.
No native triangulated realization of this example is asserted.
\end{proposition}

\begin{proof}
Appendix~\ref{app:triangle-data} and the supplied file
\texttt{triangle\_tree\_sensitivity.json} give the
complete graph, all triangle-column supports, and two explicit spanning
trees.  Each tree has eight edges.  For the successful tree, the file
lists 17 unit-column peels removing all 17 surviving rows.

For the second tree, the recorded eight peels leave the nonempty row set
\[
                         \{0,4,6,7,13,17,21,22,23\}.
\]
Its column totals, in the supplied column order, are
\[
             (2,3,0,0,2,0,2,2,0,0,2,0,0,2,0,2,2).
\]
No total is one, so this is a stopping set.  The stopping-set persistence
theorem of \cite{IncomingDualNormalization} excludes every complete
positive raw cancellation schedule, including schedules with positive
fill.  The failure is therefore independent of the particular displayed
peel order.

All finite assertions are checked by the supplied standard-library
script \path{check_triangle_witness.py}, which imports no selector,
peeling backend, or topology library.  It checks that every column
consists of three distinct edges with degree two at exactly three graph
vertices, that each proposed tree is connected and has eight edges,
that every recorded pivot has live column total one, and that the final
residual has the displayed stopping totals.  A triangle has even degree
at every vertex and exactly three occurrences, establishing both claimed
necessary compatibility conditions.
\end{proof}

The witness is an explicit finite counterexample to an abstract
invariance claim.  It does not establish how often dual choice matters
on geometrically authenticated sources after other normalization
improvements, nor does it exclude a stronger invariance theorem for a
more restricted native family.  Its role is to keep the distinction
between the proved local parity constraints and a full geometric
realization theorem explicit.
